\documentclass[10pt]{beamer}
\usetheme{metropolis}
\usepackage{appendixnumberbeamer}
\usepackage{booktabs}
\usepackage[scale=2]{ccicons}
\usepackage{pgfplots}
\usepgfplotslibrary{dateplot}
\usepackage{xspace}
\newcommand{\themename}{\textbf{\textsc{metropolis}}\xspace}
\usepackage{caption}
\usepackage{caption}

% --- Custom Header for UCSD Tailoring ---
\title{Behavioral Discounting of Delayed Protection: Evidence from a
Pandemic Vaccination Discrete Choice Experiment}

\subtitle{Job Market Paper --- Postdoctoral Candidate Presentation \\
\vspace{0.2em}
\small Herbert Wertheim School of Public Health | UC San Diego}

\author{Yichao Jin}
\institute{The University of Texas at Dallas}
\date{\today}

\begin{document}

\maketitle

% -----------------------------------------------------------------------------
% SECTION I: INTRODUCTION
% -----------------------------------------------------------------------------
\section{Introduction}

\begin{frame}{Motivation: The Intertemporal Dimension of Health}
    \begin{itemize}
        \item Global vaccination research: Focused on "Hesitancy" vs. "Acceptance."
        \item \textbf{Our Contribution}: Focus on the \alert{Intertemporal Dimension}.
        \item Does the \textit{waiting time} to protection significantly discount the perceived value of vaccination?
        \item This is critical for logistics, booster scheduling, and managing future outbreaks.
    \end{itemize}
\end{frame}

\begin{frame}{Research Questions}
    \begin{itemize}
        \item \textbf{Structural Identification}: Can we recover structural intertemporal parameters ($\delta, \kappa$) specifically for vaccination benefits?
        \item \textbf{Behavioral Pattern}: Is vaccination discounting better characterized by \alert{constant (Exponential)} or \alert{present-biased (Hyperbolic)} preferences?
        \item \textbf{Heterogeneity}: How are these structural preferences distributed across different population subgroups?
    \end{itemize}
\end{frame}

% -----------------------------------------------------------------------------
% SECTION II: METHODOLOGY
% -----------------------------------------------------------------------------
% -----------------------------------------------------------------------------
% SECTION II: METHODOLOGY: THE BEHAVIORAL TOOLBOX
% -----------------------------------------------------------------------------
\section{Methodology: The Behavioral Toolbox}

\begin{frame}{Target Population and Study Setting}
\begin{itemize}
\item \textbf{The Epicenter Context (Wuhan)}:
\begin{itemize}
\item Unique historical proximity to the initial COVID-19 outbreak.
\item High \alert{risk-salience} among residents provides a robust environment for eliciting intertemporal preferences.
\end{itemize}
\item \textbf{Strategic Sampling Window (Early 2021)}:
\begin{itemize}
\item Captured preferences during the critical voluntary phase of vaccination.
\item Reflects \textit{latent preferences} rather than responses to administrative mandates.
\end{itemize}
\end{itemize}
\end{frame}

\begin{frame}{Survey Design and Structure}
\begin{itemize}
\item \textbf{Platform}: Online survey panel in Wuhan ($N=1,027$).
\item \textbf{Instrument Structure}:
\begin{enumerate}
\item \textbf{Socio-Demographics}: Age, gender, income, education, residence.
\item \textbf{Health \& Attitudes}: Biological vulnerability, institutional trust, preventive-health routines.
\item \textbf{The DCE}: 6 choice tasks per respondent (drawn from a D-efficient design).
\item \textbf{Quality Control}: Attention filters, time-to-complete, and dominance checks.
\end{enumerate}
\end{itemize}
\end{frame}

\begin{frame}{Attribute and Level Development}
\begin{itemize}
\item \textbf{Stage 1: Evidence Synthesis}: Review of 50+ papers and policy reports on vaccine hesitancy and time preferences.
\item \textbf{Stage 2: Expert Consultation}: Qualitative refinement with public health practitioners to ensure policy relevance.
\item \textbf{Stage 3: Pilot Study ($n=60$)}:
\begin{itemize}
\item Cognitive interviews to verify comprehension of "Time to Protection."
\item Used pilot data to elicit \alert{priors} for the final experimental design.
\end{itemize}
\end{itemize}
\end{frame}

\begin{frame}{Experimental Design: Attributes and Levels}
\begin{table}
\centering
\small
\begin{tabular}{ll}
\toprule
Attribute & Levels \\
\midrule
\textbf{Vaccination Delay} ($T$) & \alert{0, 1, 3, 6 months} \\
Vaccine Efficacy ($E$) & 50\%, 70\%, 95\% \\
Side Effects ($S$) & Mild, Moderate, Severe \\
Vaccine Origin ($O$) & Domestic, Imported \\
Incentives ($M$) & 50, 200, 800 RMB \\
\bottomrule
\end{tabular}
\caption{Final attributes justified by policy roll-out constraints.}
\end{table}
\end{frame}

\begin{frame}{DCE Construction: Sample Choice Task}
\begin{table}
\centering
\small
\begin{tabular}{|l|c|c|c|}
\hline
\textbf{Attribute} & \textbf{Vaccine A} & \textbf{Vaccine B} & \textbf{Opt-out} \\
\hline
Waiting Time & 0 months & 3 months & \\
Efficacy & 95\% & 95\% & \\
Side Effects & Mild & Mild & \textbf{I would} \\
Origin & Domestic & Imported & \textbf{choose} \\
Incentive & 50 RMB & 200 RMB & \textbf{neither.} \\
\hline
\textbf{Your Choice?} & $\square$ & $\square$ & $\square$ \\
\hline
\end{tabular}
\end{table}
\begin{itemize}
\item \textbf{The Task}: Choose the preferred vaccine profile or opt-out.
\item \textbf{Design Efficiency}: D-efficient design using pilot priors.
\end{itemize}
\end{frame}

\begin{frame}{Variables and Econometric Specification}
\begin{itemize}
\item \textbf{Dependent Variable ($Y_{nit}$)}:
\begin{itemize}
\item Discrete choice for alternative $j$ in set $t$ by respondent $n$.
\end{itemize}
\item \textbf{Independent Variables (Attributes)}:
\begin{itemize}
\item \textbf{Delay}: Treated as a continuous variable for structural mapping.
\item \textbf{Money}: Marginal utility of incentives used as the MWTA denominator.
\item \textbf{Non-linearities}: Dummies for Efficacy and Side Effects.
\end{itemize}

\item \textbf{Mixed Logit (MXL)}: $\beta_n \sim f(\beta|\theta)$ using 1,000 Halton draws to capture random preference variation.
\item \textbf{Identification}: Structural recovery of ($\delta, \kappa$) via non-linear mapping of the delay coefficient.
\end{itemize}
\end{frame}
% -----------------------------------------------------------------------------
% SECTION III: RESULTS (JMP)
% -----------------------------------------------------------------------------
% --- Section III: Results I: Completed Study (JMP) ---
\section{Results I: Behavioral Discounting Findings}

\begin{frame}{Model Performance: Comparative Fit}
    \begin{itemize}
        \item \textbf{Question:} Which discounting functional form best captures vaccine timing decisions?
        \item \textbf{Metric:} Log-Likelihood (LL) and AIC/BIC criteria.
    \end{itemize}
    \begin{table}[]
        \centering
        \begin{tabular}{lccc}
            \hline
            \textbf{Model} & \textbf{LL} & \textbf{AIC} & \textbf{Log-Rank} \\ \hline
            Exponential ($V = \delta^D \cdot U$) & -4,521.2 & 9,046.4 & -- \\
            \textbf{Hyperbolic ($V = \frac{1}{1+\kappa D} \cdot U$)} & \textbf{-4,288.5} & \textbf{8,581.0} & \textbf{*** (p<0.01)} \\ \hline
        \end{tabular}
    \end{table}
    \begin{itemize}
        \item \textbf{Result:} Hyperbolic discounting significantly outperforms Exponential models.
        \item \textbf{Implication:} Presence of "Present Bias" in health protection delay.
    \end{itemize}
\end{frame}

\begin{frame}{Structural Estimates: Preference Parameters}
    \begin{itemize}
        \item \textbf{Model:} $U_{inpt} = \alpha \cdot \text{Effectiveness}_n + \beta \cdot \text{Risk}_p + \text{Discounting}(\kappa, \delta) \dots$
    \end{itemize}
    \begin{columns}
        \begin{column}{0.45\textwidth}
            \begin{block}{Parameter Estimates}
                \begin{itemize}
                    \item $\hat{\delta}$ (Base Discount): 0.92 [0.89, 0.95]
                    \item $\hat{\kappa}$ (Hyperbolic Rate): 0.14 [0.11, 0.17]
                \end{itemize}
            \end{block}
            \begin{itemize}
                \item \footnotesize{Note: Estimates derived from Mixed Logit specification to account for unobserved heterogeneity.}
            \end{itemize}
        \end{column}
        \begin{column}{0.55\textwidth}
            \centering
            % 自动匹配你目录下的参数分布/异质性图表
            \includegraphics[width=\textwidth]{Subgroup Variation in Hyperbolic Discount Parameters.png}
            \\ \vspace{0.2cm}
            \captionof{figure}{\footnotesize{Distribution of $\kappa$ across sub-populations}}
        \end{column}
    \end{columns}
\end{frame}

\begin{frame}{The "Behavioral Delay Multiplier" (BDM) Reveal}
    \begin{itemize}
        \item \textbf{Definition:} The psychological weight assigned to 1 unit of biological delay.
        \item \textbf{Core Finding:} \color{red}\textbf{BDM = 1.29}\color{black}
    \end{itemize}
    \begin{exampleblock}{Interpretation}
        An average individual perceives a 14-day delay in antibody production as equivalent to an \textbf{18.1-day} (14 $\times$ 1.29) delay in utility. 
    \end{exampleblock}
    \begin{itemize}
        \item This 29\% "Psychological Markup" explains the observed hesitancy even for high-efficacy vaccines.
    \end{itemize}
\end{frame}

\begin{frame}{Scientific Significance: Why Structural Modeling?}
    \begin{itemize}
        \item \textbf{Beyond Correlation:} Reduced-form models only show "delay matters." 
        \item \textbf{The Unique Contribution:}
        \begin{enumerate}
            \item \textbf{Recovering Latent Parameters:} We isolate biological wait-time preferences from general risk aversion.
            \item \textbf{Policy Simulation:} Allows predicting the impact of "Express Vaccines" vs. "High Efficacy" on uptake.
            \item \textbf{Counterfactuals:} What happens if BDM is reduced via information intervention?
        \end{enumerate}
    \end{itemize}
\end{frame}

\begin{frame}{Mechanism: The Buffer Effect of Institutional Trust}
    \begin{itemize}
        \item Does trust mitigate the psychological cost of waiting?
    \end{itemize}
    \begin{table}[]
        \centering
        \begin{tabular}{lc}
            \hline
            \textbf{Sub-group} & \textbf{Behavioral Delay Multiplier (BDM)} \\ \hline
            High Institutional Trust & \textbf{1.19} (Low psychological markup) \\
            Low Institutional Trust & \textbf{1.55} (Exaggerated wait-time cost) \\ \hline
        \end{tabular}
    \end{table}
    \begin{itemize}
% 
\item \textbf{Insight:} Trust acts as a "behavioral stabilizer," smoothing the discounting of future protection.
    \end{itemize}
\end{frame}

\begin{frame}{Visualizing Evidence: Iso-acceptance Heatmap}
    \centering
    % 插入等验收热力图
    \includegraphics[width=0.85\textwidth]{Iso-acceptance heatmap.png}
    
    \vspace{0.3cm}
    \begin{itemize}
        \item \textbf{Iso-acceptance Lines:} Represent combinations of Effectiveness and Delay that yield the same probability of uptake.
        \item \textbf{The Trade-off:} The steepness of the curves quantifies the \textbf{Marginal Rate of Substitution (MRS)} between performance and timing.
        \item \textbf{Key Finding:} To compensate for a \textbf{1-week increase} in protection delay, vaccine effectiveness must increase by approximately \textbf{2.5\%} to maintain the same level of acceptance.
    \end{itemize}
   \end{frame}

\begin{frame}{Validity Diagnostics I: Theoretical Consistency}
    \begin{itemize}
        \item \textbf{Theoretical Validity (Sign Check):} 
        \begin{itemize}
            \item Marginal utilities follow expected directions: $\beta_{Effectiveness} > 0$, $\beta_{Incentive} > 0$, and $\beta_{Delay} < 0$.
        \end{itemize}
        \item \textbf{Internal Consistency (Dominance Test):}
        \begin{itemize}
            \item A "dominant" choice set was included (higher efficacy + lower delay + higher incentive).
            \item \textbf{92\%} of respondents passed the dominance test, indicating high comprehension and engagement.
        \end{itemize}
    \end{itemize}
    \begin{columns}
        \begin{column}{0.5\textwidth}
             \begin{itemize}
                \item \textbf{Test-Retest Reliability:}
                \begin{itemize}
                    \item 85.5\% consistency rate on repeated choice tasks.
                \end{itemize}
            \end{itemize}
        \end{column}
        \begin{column}{0.5\textwidth}
            \centering
            \includegraphics[width=\textwidth]{correlation among Individual variables4.png}
            \caption{\footnotesize{External Validity: Correlation between estimated preferences and individual characteristics.}}
        \end{column}
    \end{columns}
\end{frame}

\begin{frame}{Validity Diagnostics II: Attribute Processing}
    \begin{itemize}
        \item \textbf{Attribute Attendance:} 
        \begin{itemize}
            \item Self-reported importance rankings (stated) align with estimated marginal utilities (revealed). 
            \item Cash and Delay were ranked as the top-two decision drivers.
        \end{itemize}
    \end{itemize}
    \begin{columns}
        \begin{column}{0.55\textwidth}
            \centering
            \includegraphics[width=\textwidth]{Subgroup differences in time sensitivity.png}
            \caption{\footnotesize{Stability Check: Parameter estimates remain robust across different sub-samples.}}
        \end{column}
        \begin{column}{0.45\textwidth}
            \begin{block}{Model Robustness}
                \begin{itemize}
                    \item \textbf{Non-linearity:} Inclusion of squared terms for delay confirms the hyperbolic decay pattern.
                    \item \textbf{Scale Heterogeneity:} Model fit improves when accounting for task-level variance.
                \end{itemize}
            \end{block}
        \end{column}
    \end{columns}
\end{frame}
% SECTION IV: POLICY IMPACT EVALUATION
% -----------------------------------------------------------------------------
\section{Results II: Policy Impact Evaluation}

\begin{frame}{Welfare Analysis: The Shadow Price of Waiting}
    \begin{itemize}
        \item \textbf{Shadow Price}: Translating time to money via MWTA.
        \item \textbf{Mean MWTA = 111.41 RMB/month}.
        \item Quantifies the welfare loss of logistical inefficiency across different socio-economic strata.
    \end{itemize}
\end{frame}

\begin{frame}{Counterfactual Simulations: Speed vs. Incentives}
    \begin{itemize}
        \item Evaluating the efficiency of \alert{logistical speed-ups} vs. \alert{monetary compensation}.
        \item Framework for optimal budget allocation under constraints.
    \end{itemize}
    \begin{center}
        \includegraphics[width=0.55\textwidth]{essay3_equity_budget_allocation.png}
    \end{center}
\end{frame}

% -----------------------------------------------------------------------------
% SECTION V: STRATEGIC BRIDGE (UCSD SYNERGY)
% -----------------------------------------------------------------------------
\section{Future Research Agenda: Bridging Behavioral Models to Drug Policy}

\begin{frame}{Transferability I: Present-Bias in Addiction}
    \begin{itemize}
        \item \textbf{Fundamental Link}: The intertemporal conflict in vaccination is identical to choices in substance use cessation.
        \item \textbf{Modeling Addiction}: Applying $\hat{\kappa}$ and the BDM to evaluate "冲动性决策" (impulsivity) in tobacco and cannabis use.
        \item \textit{Value}: Identifying the optimal intervention window for at-risk populations.
    \end{itemize}
\end{frame}

\begin{frame}{Transferability II: Regulatory Barriers & Shadow Price}
    \begin{itemize}
        \item \textbf{Access as a Cost}: Regulatory hurdles (age verification, tax increases, flavoring bans) act as a "time tax" on users.
        \item \textbf{Impact Evaluation}: Using the MWTA framework to predict responses to new drug regulations across socio-economic subgroups.
    \end{itemize}
\end{frame}

\section{Conclusion}

\begin{frame}{Conclusion and Team Fit}
    \begin{itemize}
        \item \textbf{Contributions}: Structural evidence of hyperbolic discounting; Original policy metrics (BDM/MWTA).
        \item \textbf{UCSD Fit}: Eager to apply this behavioral toolkit to the Shi Lab's impactful drug policy research.
    \end{itemize}
    \begin{center}
        \Huge Thank you! \\
        \vspace{1em}
        \large Yichao Jin \\
        \small Yichao.Jin@UTDallas.edu
    \end{center}
\end{frame}

\end{document}
